
\documentclass[11pt,a4paper]{article}

% ============================================================
% PACKAGES
% ============================================================

\usepackage[utf8]{inputenc}
\usepackage[T1]{fontenc}
\usepackage{lmodern}

\usepackage{amsmath}
\usepackage{amssymb}
\usepackage{bm}

\usepackage{geometry}
\geometry{
    a4paper,
    left=25mm,
    right=25mm,
    top=25mm,
    bottom=25mm
}

\usepackage{booktabs}
\usepackage{array}
\usepackage{longtable}
\usepackage{float}
\usepackage{siunitx}

\usepackage{graphicx}
\usepackage{caption}
\usepackage{subcaption}

\usepackage{xcolor}
\usepackage{hyperref}
\usepackage{microtype}

\usepackage{enumitem}
\usepackage{fancyhdr}

% ============================================================
% DOCUMENT SETTINGS
% ============================================================

\hypersetup{
    colorlinks=true,
    linkcolor=black,
    citecolor=black,
    urlcolor=blue,
    pdftitle={AVL Dynamic Stability and Eigenmode Analysis},
    pdfauthor={Aircraft Stability Analysis}
}

\setlength{\parindent}{0pt}
\setlength{\parskip}{6pt}

\setlist[itemize]{leftmargin=1.5em}
\setlist[enumerate]{leftmargin=1.5em}

\sisetup{
    scientific-notation = true,
    round-mode = places,
    round-precision = 4
}

% ============================================================
% HEADER / FOOTER
% ============================================================

\pagestyle{fancy}
\fancyhf{}
\lhead{Dynamic Stability Analysis}
\rhead{AVL}
\cfoot{\thepage}

% ============================================================
% DOCUMENT
% ============================================================

\begin{document}


% ============================================================
\title{Dynamic Stability Analysis Using AVL:}
% ============================================================


The main objectives of this analysis are:

\begin{enumerate}
    \item To obtain the aerodynamic stability derivatives of the aircraft.
    \item To determine the longitudinal and lateral-directional stability
    characteristics.
    \item To determine the linearized state-space representation of the
    aircraft.
    \item To calculate the eigenvalues of the linearized dynamic system.
    \item To identify the aircraft's principal dynamic modes.
    \item To evaluate the damping, natural frequency, and stability of each
    mode.
\end{enumerate}

No control-surface actuation or feedback control system was introduced in
this analysis. Therefore, the calculated eigenvalues represent the
uncontrolled aircraft dynamics at the selected operating condition.

% ============================================================
\section{Aircraft Configuration and Reference Data}
% ============================================================

The AVL model represents the complete aerodynamic configuration of the
aircraft, consisting of the main wing, horizontal tail, and vertical tail.

The aerodynamic reference quantities used in the AVL model are summarized
in Table~\ref{tab:reference_data}.

\begin{table}[H]
\centering
\caption{Aircraft reference quantities used in AVL.}
\label{tab:reference_data}
\begin{tabular}{lc}
\toprule
Parameter & Value \\
\midrule
Reference area, $S_{\mathrm{ref}}$
    & $0.08001~\mathrm{m^2}$ \\

Reference chord, $c_{\mathrm{ref}}$
    & $0.12969~\mathrm{m}$ \\

Reference span, $b_{\mathrm{ref}}$
    & $0.63000~\mathrm{m}$ \\

$X_{\mathrm{ref}}$
    & $0.07879~\mathrm{m}$ \\

$Y_{\mathrm{ref}}$
    & $0.00000~\mathrm{m}$ \\

$Z_{\mathrm{ref}}$
    & $0.00697~\mathrm{m}$ \\

Air density, $\rho$
    & $1.225~\mathrm{kg/m^3}$ \\

Gravity, $g$
    & $9.81~\mathrm{m/s^2}$ \\

Aircraft mass, $m$
    & $0.514~\mathrm{kg}$ \\

\bottomrule
\end{tabular}
\end{table}

The total mass was obtained from the complete mass breakdown:

\begin{equation}
\begin{aligned}
m &=
0.125+0.018+0.010+0.134\\
&\quad+0.075+0.028+0.124\\
&=0.514~\mathrm{kg}.
\end{aligned}
\end{equation}

The calculated center of gravity is approximately

\begin{equation}
X_{CG}=0.07880~\mathrm{m},
\end{equation}

\begin{equation}
Y_{CG}\approx0,
\end{equation}

and

\begin{equation}
Z_{CG}=0.006975~\mathrm{m}.
\end{equation}

% ============================================================
\section{Mass and Inertia Properties}
% ============================================================

The mass and inertia properties supplied to AVL were obtained from the
aircraft mass model.

The principal inertia values used in the dynamic analysis were

\begin{align}
I_{xx} &= 0.003232~\mathrm{kg\,m^2},\\
I_{yy} &= 0.004960~\mathrm{kg\,m^2},\\
I_{zz} &= 0.008001~\mathrm{kg\,m^2}.
\end{align}

The corresponding inertia tensor was

\begin{equation}
\mathbf{I}
=
\begin{bmatrix}
0.003232 & 0 & 1.89\times10^{-5}\\
0 & 0.004960 & 0\\
1.89\times10^{-5} & 0 & 0.008001
\end{bmatrix}
\mathrm{kg\,m^2}.
\end{equation}

These values were used by AVL in constructing the linearized aircraft
equations of motion.

% ============================================================
\section{AVL Analysis Setup}
% ============================================================

The aircraft was analyzed using AVL version 3.37.

The aerodynamic model consisted of:

\begin{itemize}
    \item Main wing,
    \item Horizontal tail,
    \item Vertical tail.
\end{itemize}

The model contained 83 aerodynamic strips and 809 vortex elements.

The following operating condition was selected for the dynamic analysis:

\begin{equation}
V=20.0~\mathrm{m/s},
\end{equation}

with

\begin{equation}
\rho=1.225~\mathrm{kg/m^3}.
\end{equation}

The corresponding Run Case 1 solution produced a required lift coefficient
of

\begin{equation}
C_L=0.25723.
\end{equation}

% ============================================================
\section{Aerodynamic Operating Point}
% ============================================================

At the selected operating condition, AVL returned the following
aerodynamic solution:

\begin{table}[H]
\centering
\caption{Aerodynamic operating point obtained from AVL.}
\label{tab:operating_point}
\begin{tabular}{lc}
\toprule
Quantity & Value \\
\midrule
Velocity, $V$ & $20.0~\mathrm{m/s}$ \\
Angle of attack, $\alpha$ & $-8.303^\circ$ \\
Lift coefficient, $C_L$ & $0.25723$ \\
Drag coefficient, $C_D$ & $0.01696$ \\
Pitching moment coefficient, $C_m$ & $+0.21928$ \\
$C_X$ & $-0.05393$ \\
$C_Z$ & $-0.25208$ \\
Induced drag coefficient, $C_{D_i}$ & $0.016958$ \\
Span efficiency, $e$ & $0.2679$ \\
\bottomrule
\end{tabular}
\end{table}

% ============================================================
\subsection{Operating Point Limitation}
% ============================================================

An important distinction must be made between the lift condition imposed
during the AVL analysis and a complete longitudinal trim condition.

The calculated pitching moment coefficient was

\begin{equation}
C_m=+0.21928.
\end{equation}

Since

\begin{equation}
C_m\neq0,
\end{equation}

the selected operating point is not a complete longitudinal pitching
moment equilibrium.

Therefore, the eigenvalues presented in this report should be interpreted
as the linearized dynamic characteristics of the aircraft around the
selected AVL operating point.

This limitation is explicitly stated because the purpose of the present
analysis is to extract and evaluate the aircraft's dynamic modes, while
the final trimmed aerodynamic condition is investigated separately.

% ============================================================
\section{Static Stability Derivatives}
% ============================================================

AVL was used to calculate the aircraft stability derivatives.

The resulting derivatives are presented in
Table~\ref{tab:static_derivatives}.

\begin{table}[H]
\centering
\caption{Static stability derivatives obtained from AVL.}
\label{tab:static_derivatives}
\begin{tabular}{lr}
\toprule
Derivative & Value \\
\midrule
$C_{L_\alpha}$ & $4.757761$ \\
$C_{L_\beta}$ & $0$ \\
$C_{Y_\beta}$ & $-0.246263$ \\
$C_{l_\beta}$ & $-0.081797$ \\
$C_{m_\alpha}$ & $-0.895599$ \\
$C_{m_\beta}$ & $0$ \\
$C_{n_\beta}$ & $0.112811$ \\
\bottomrule
\end{tabular}
\end{table}

The negative pitching moment derivative

\begin{equation}
C_{m_\alpha}=-0.895599
\end{equation}

indicates positive longitudinal static stability.

The lateral-directional derivatives also indicate favorable static
stability characteristics:

\begin{equation}
C_{l_\beta}=-0.081797,
\end{equation}

and

\begin{equation}
C_{n_\beta}=+0.112811.
\end{equation}

These results indicate restoring rolling and yawing moments following
small sideslip disturbances.

% ============================================================
\section{Dynamic Stability Derivatives}
% ============================================================

The dynamic stability derivatives calculated by AVL are shown in
Table~\ref{tab:dynamic_derivatives}.

\begin{table}[H]
\centering
\caption{Dynamic stability derivatives obtained from AVL.}
\label{tab:dynamic_derivatives}
\begin{tabular}{lr}
\toprule
Derivative & Value \\
\midrule
$C_{L_p}$ & $0$ \\
$C_{L_q}$ & $9.007163$ \\
$C_{L_r}$ & $0$ \\
$C_{Y_p}$ & $-0.236351$ \\
$C_{Y_q}$ & $-0.000002$ \\
$C_{Y_r}$ & $0.309552$ \\
$C_{l_p}$ & $-0.426048$ \\
$C_{l_q}$ & $0$ \\
$C_{l_r}$ & $0.136908$ \\
$C_{m_p}$ & $0$ \\
$C_{m_q}$ & $-16.254972$ \\
$C_{m_r}$ & $0$ \\
$C_{n_p}$ & $0.016380$ \\
$C_{n_q}$ & $-0.000002$ \\
$C_{n_r}$ & $-0.155319$ \\
\bottomrule
\end{tabular}
\end{table}

The roll damping derivative

\begin{equation}
C_{l_p}=-0.426048
\end{equation}

indicates a restoring aerodynamic moment opposing roll rate.

Similarly, the pitch damping derivative

\begin{equation}
C_{m_q}=-16.254972
\end{equation}

indicates strong aerodynamic damping in pitch.

The yaw damping derivative is

\begin{equation}
C_{n_r}=-0.155319,
\end{equation}

which also contributes to damping of the lateral-directional response.

% ============================================================
\section{Neutral Point and Static Margin}
% ============================================================

The neutral point calculated by AVL was

\begin{equation}
X_{NP}=0.103210~\mathrm{m}.
\end{equation}

The aircraft center of gravity was

\begin{equation}
X_{CG}=0.07880~\mathrm{m}.
\end{equation}

The static margin is therefore

\begin{equation}
SM=
\frac{X_{NP}-X_{CG}}{c_{\mathrm{ref}}}.
\end{equation}

Substituting the measured values gives

\begin{equation}
SM=
\frac{0.103210-0.07880}{0.12969}
\approx0.188.
\end{equation}

Thus,

\begin{equation}
\boxed{SM\approx18.8\%}
\end{equation}

which indicates a positive static margin.

This result agrees with the negative value of $C_{m_\alpha}$ and confirms
positive longitudinal static stability.

% ============================================================
\section{Linearized State-Space Model}
% ============================================================

AVL generated the linearized equations of motion in state-space form:

\begin{equation}
\dot{\mathbf{x}}=\mathbf{A}\mathbf{x}.
\end{equation}

The state vector used by AVL was

\begin{equation}
\mathbf{x}
=
\begin{bmatrix}
u & w & q & \theta &
v & p & r & \phi &
x & y & z & \psi
\end{bmatrix}^{T}.
\end{equation}

The system matrix obtained directly from AVL was

\begin{equation}
\mathbf{A}
=
\begin{bmatrix}
-0.3538 & -0.9998 & 2.7917 & -9.8100 & 0 & 0 & 0 & 0 & 0 & 0 & 0 & 0\\
-2.2079 & -8.6195 & 18.2728 & 0 & 0 & 0 & 0 & 0 & 0 & 0 & 0 & 0\\
7.1011 & -22.1349 & -24.5711 & 0 & 0 & 0 & 0 & 0 & 0 & 0 & 0 & 0\\
0 & 0 & 1.0000 & 0 & 0 & 0 & 0 & 0 & 0 & 0 & 0 & 0\\
0 & 0 & 0 & 0 & -0.4689 & -3.0633 & -19.5583 & 9.8100 & 0 & 0 & 0 & 0\\
0 & 0 & 0 & 0 & -10.9110 & -22.1584 & 9.3661 & 0 & 0 & 0 & 0\\
0 & 0 & 0 & 0 & 9.1177 & 0.6975 & -4.1497 & 0 & 0 & 0 & 0\\
0 & 0 & 0 & 0 & 0 & 1.0000 & 0 & 0 & 0 & 0 & 0 & 0\\
1.0000 & 0 & 0 & -2.8883 & 0 & 0 & 0 & 0 & 0 & 0 & 0 & 0\\
0 & 0 & 0 & 0 & 1.0000 & 0 & 0 & 2.8883 & 0 & 0 & 0 & 19.7903\\
0 & 1.0000 & 0 & -19.7903 & 0 & 0 & 0 & 0 & 0 & 0 & 0 & 0\\
0 & 0 & 0 & 0 & 0 & 0 & 1.0000 & 0 & 0 & 0 & 0 & 0
\end{bmatrix}.
\end{equation}

The eigenvalues of this matrix provide the natural dynamic modes of the
aircraft.

% ============================================================
\section{Eigenvalue Analysis}
% ============================================================

For a linearized dynamic system, each eigenvalue may be written as

\begin{equation}
\lambda=\sigma\pm i\omega,
\end{equation}

where $\sigma$ is the real part and $\omega$ is the imaginary part.

The sign of $\sigma$ determines stability:

\begin{equation}
\sigma<0
\quad\Rightarrow\quad
\text{stable mode},
\end{equation}

while

\begin{equation}
\sigma>0
\quad\Rightarrow\quad
\text{unstable mode}.
\end{equation}

The eigenvalues obtained directly from AVL are shown in
Table~\ref{tab:eigenvalues}.

\begin{table}[H]
\centering
\caption{Eigenvalues obtained from AVL Run Case 1.}
\label{tab:eigenvalues}
\begin{tabular}{ccl}
\toprule
Mode & Eigenvalue & Type \\
\midrule
1 &
$-23.800234$ &
Real \\

2--3 &
$-1.5364187\pm13.028574i$ &
Complex pair \\

4--5 &
$-16.775574\pm17.971889i$ &
Complex pair \\

6 &
$+0.096084811$ &
Real \\

7--8 &
$+0.0033943544\pm1.3366623i$ &
Complex pair \\

\bottomrule
\end{tabular}
\end{table}

% ============================================================
\section{Dynamic Mode Identification}
% ============================================================

The eigenvalue spectrum corresponds to five principal aircraft dynamic
modes:

\begin{enumerate}
    \item Roll subsidence,
    \item Dutch roll,
    \item Short period,
    \item Spiral mode,
    \item Phugoid.
\end{enumerate}

Complex conjugate eigenvalues correspond to oscillatory modes, whereas
real eigenvalues correspond to non-oscillatory modes.

% ============================================================
\subsection{Roll Subsidence}
% ============================================================

The first eigenvalue is

\begin{equation}
\lambda_{\mathrm{roll}}
=
-23.800234.
\end{equation}

The negative real eigenvalue indicates a stable, non-oscillatory roll
subsidence mode.

The characteristic time constant is

\begin{equation}
\tau_{\mathrm{roll}}
=
\frac{1}{|\sigma|}
=
\frac{1}{23.800234}
\approx0.0420~\mathrm{s}.
\end{equation}

The corresponding half-life is

\begin{equation}
t_{1/2}
=
\frac{\ln2}{23.800234}
\approx0.0291~\mathrm{s}.
\end{equation}

Therefore, the roll disturbance decays very rapidly.

% ============================================================
\subsection{Dutch Roll}
% ============================================================

The Dutch-roll eigenvalue pair is

\begin{equation}
\lambda_{\mathrm{DR}}
=
-1.5364187
\pm13.028574i.
\end{equation}

The natural frequency is calculated from

\begin{equation}
\omega_n
=
\sqrt{\sigma^2+\omega^2}.
\end{equation}

Therefore,

\begin{equation}
\omega_n
=
\sqrt{(-1.5364187)^2+(13.028574)^2}
\approx13.119~\mathrm{rad/s}.
\end{equation}

The damping ratio is

\begin{equation}
\zeta
=
-\frac{\sigma}{\omega_n}
\approx0.117.
\end{equation}

The oscillation period is

\begin{equation}
T=
\frac{2\pi}{|\omega|}
=
\frac{2\pi}{13.028574}
\approx0.482~\mathrm{s}.
\end{equation}

The Dutch-roll mode is therefore stable but lightly damped.

% ============================================================
\subsection{Short-Period Mode}
% ============================================================

The short-period eigenvalue pair is

\begin{equation}
\lambda_{\mathrm{SP}}
=
-16.775574
\pm17.971889i.
\end{equation}

The natural frequency is

\begin{equation}
\omega_n
=
\sqrt{(-16.775574)^2+(17.971889)^2}
\approx24.583~\mathrm{rad/s}.
\end{equation}

The damping ratio is

\begin{equation}
\zeta
=
-\frac{-16.775574}{24.583}
\approx0.682.
\end{equation}

The oscillation period is

\begin{equation}
T=
\frac{2\pi}{17.971889}
\approx0.350~\mathrm{s}.
\end{equation}

The short-period mode is therefore stable and strongly damped.

% ============================================================
\subsection{Spiral Mode}
% ============================================================

The sixth eigenvalue is

\begin{equation}
\lambda_{\mathrm{spiral}}
=
+0.096084811.
\end{equation}

Because the real part is positive, the spiral mode is unstable.

The disturbance amplitude behaves approximately as

\begin{equation}
x(t)\propto
e^{0.096084811t}.
\end{equation}

The corresponding doubling time is

\begin{equation}
t_2
=
\frac{\ln2}{0.096084811}
\approx7.22~\mathrm{s}.
\end{equation}

Thus, the amplitude associated with the spiral mode doubles approximately
every $7.22$ seconds in the absence of corrective control input.

This is the most significant dynamic instability identified in the
present analysis.

% ============================================================
\subsection{Phugoid}
% ============================================================

The phugoid eigenvalue pair is

\begin{equation}
\lambda_{\mathrm{PH}}
=
+0.0033943544
\pm1.3366623i.
\end{equation}

The natural frequency is

\begin{equation}
\omega_n
=
\sqrt{(0.0033943544)^2+(1.3366623)^2}
\approx1.33666~\mathrm{rad/s}.
\end{equation}

The damping ratio is

\begin{equation}
\zeta
=
-\frac{0.0033943544}{1.33666}
\approx-0.00254.
\end{equation}

The negative damping ratio indicates that the mode is very slightly
unstable.

The oscillation period is

\begin{equation}
T=
\frac{2\pi}{1.3366623}
\approx4.70~\mathrm{s}.
\end{equation}

The amplitude doubling time is

\begin{equation}
t_2
=
\frac{\ln2}{0.0033943544}
\approx204.2~\mathrm{s}.
\end{equation}

Therefore, although the phugoid is mathematically unstable, its growth
rate is extremely small compared with that of the spiral mode.

% ============================================================
\section{Dynamic Mode Summary}
% ============================================================

The principal dynamic characteristics are summarized in
Table~\ref{tab:mode_summary}.

\begin{table}[H]
\centering
\caption{Summary of the aircraft dynamic modes.}
\label{tab:mode_summary}
\begin{tabular}{lccccc}
\toprule
Mode & $\sigma$ & $\omega$ & $\omega_n$ & $\zeta$ & Stability \\
\midrule
Roll subsidence
& $-23.800$
& $0$
& $23.800$
& $1.000$
& Stable \\

Dutch roll
& $-1.536$
& $\pm13.029$
& $13.119$
& $0.117$
& Stable \\

Short period
& $-16.776$
& $\pm17.972$
& $24.583$
& $0.682$
& Stable \\

Spiral
& $+0.0961$
& $0$
& $0.0961$
& --
& Unstable \\

Phugoid
& $+0.00339$
& $\pm1.337$
& $1.337$
& $-0.00254$
& Very slightly unstable \\

\bottomrule
\end{tabular}
\end{table}

% ============================================================
\section{Overall Stability Assessment}
% ============================================================

The eigenvalue analysis reveals three stable dynamic modes and two
unstable modes.

The roll subsidence mode is highly stable, with a very short time
constant. The short-period mode is also strongly damped, with a damping
ratio of approximately

\begin{equation}
\zeta_{\mathrm{SP}}\approx0.682.
\end{equation}

The Dutch-roll mode is stable but lightly damped:

\begin{equation}
\zeta_{\mathrm{DR}}\approx0.117.
\end{equation}

The principal instability is the spiral mode, which has a positive real
eigenvalue of

\begin{equation}
\sigma_{\mathrm{spiral}}
=
0.0961~\mathrm{s^{-1}}.
\end{equation}

Its doubling time of approximately $7.22$ seconds indicates a slowly
diverging lateral-directional response.

The phugoid is also mathematically unstable, but its growth rate is much
smaller:

\begin{equation}
\sigma_{\mathrm{phugoid}}
=
0.00339~\mathrm{s^{-1}}.
\end{equation}

Its corresponding doubling time is approximately $204$ seconds.
Consequently, the phugoid instability is extremely weak.

% ============================================================
\section{Comparison with Static Stability}
% ============================================================

The dynamic results are generally consistent with the static stability
characteristics.

The negative pitching-moment derivative

\begin{equation}
C_{m_\alpha}=-0.895599
\end{equation}

and the positive static margin

\begin{equation}
SM\approx18.8\%
\end{equation}

indicate positive longitudinal static stability.

This agrees with the stable short-period mode.

Similarly,

\begin{equation}
C_{l_\beta}=-0.081797
\end{equation}

and

\begin{equation}
C_{n_\beta}=+0.112811
\end{equation}

indicate favorable static lateral-directional stability.

However, the unstable spiral eigenvalue demonstrates that positive static
stability does not necessarily guarantee complete dynamic stability.

This distinction is important in aircraft design: static stability
describes the initial restoring tendency following a disturbance, whereas
dynamic stability describes the subsequent time evolution of that
disturbance.

% ============================================================
\section{Conclusions}
% ============================================================

The AVL dynamic stability analysis provided the linearized dynamic
characteristics of the aircraft at the selected operating condition.

The calculated eigenvalues were

\begin{equation}
\boxed{
\begin{aligned}
\lambda_1 &=-23.8002,\\
\lambda_{2,3}
&=-1.5364\pm13.0286i,\\
\lambda_{4,5}
&=-16.7756\pm17.9719i,\\
\lambda_6
&=+0.09608,\\
\lambda_{7,8}
&=+0.003394\pm1.33666i.
\end{aligned}}
\end{equation}

The aircraft therefore exhibits:

\begin{itemize}
    \item stable roll subsidence,
    \item stable but lightly damped Dutch roll,
    \item stable and strongly damped short-period motion,
    \item an unstable spiral mode,
    \item a very weakly unstable phugoid mode.
\end{itemize}

The spiral mode represents the primary dynamic stability concern because
it has the largest positive real eigenvalue and therefore the fastest
growth among the unstable modes.

The phugoid instability is considerably weaker and develops over a much
longer time scale.

The static stability results are favorable, with a positive static margin
of approximately $18.8\%$ and a negative $C_{m_\alpha}$.

Finally, the interpretation of the eigenvalues must account for the fact
that the selected AVL operating point satisfies the required lift
condition but does not represent a complete longitudinal trim condition,
as indicated by

\begin{equation}
C_m=+0.21928.
\end{equation}

Therefore, the results presented here should be regarded as the
linearized dynamic characteristics of the aircraft at the analyzed AVL
operating point. A final dynamic stability assessment at a fully trimmed
flight condition can be performed subsequently if required.

% ============================================================
% END OF DOCUMENT
% ============================================================

\end{document}

